% Findings ledger -- every measurement in the PointMass3D project, with its
% weight and its evidence. NOT a draft: a source document to write a draft FROM.
%
%   pdflatex findings && pdflatex findings
%
% Written against the same minimal TeX Live as paper.tex: no booktabs, no
% xcolor, no enumitem. Every optional package degrades gracefully, so this
% compiles on the laptop and on the cluster.
%
% Each entry is \label'ed as f:ID, so \ref{f:H4} works from a draft in the same
% directory tree. Cite findings by ID in notes to yourself; the IDs are stable.

\documentclass[11pt,letterpaper]{article}

\usepackage[margin=1in]{geometry}
\usepackage[utf8]{inputenc}
\usepackage{amsmath}
\usepackage{amssymb}
\usepackage{array}
\usepackage[hidelinks]{hyperref}

\IfFileExists{booktabs.sty}{\usepackage{booktabs}}{%
  \newcommand{\toprule}{\hline}
  \newcommand{\midrule}{\hline}
  \newcommand{\bottomrule}{\hline}}

% Weight tags. Small caps rather than colour, because xcolor is not installed
% here and a tag that vanishes on another machine is worse than a plain one.
\newcommand{\chip}[1]{{\footnotesize\textsc{[#1]}}}
\newcommand{\HEAD}{\chip{headline}}
\newcommand{\SUPP}{\chip{supporting}}
\newcommand{\DET}{\chip{detail}}
\newcommand{\RETR}{\chip{retracted}}
\newcommand{\OPEN}{\chip{open}}

% \finding{ID}{weight tag} followed by the claim, then optional \ev / \nt.
\newcommand{\finding}[2]{%
  \par\medskip\noindent
  \textbf{\texttt{#1}}\label{f:#1}\hspace{0.7em}#2\hspace{0.7em}\ignorespaces}
% Evidence: how it was measured, and with what.
% \raggedright + emergencystretch: the evidence lines are typewriter text full
% of long unbreakable tokens (paths, "2,582,233", "1.2e-06"), which TeX cannot
% hyphenate and will otherwise push into the margin.
\newcommand{\ev}[1]{\par\vspace{1pt}\begingroup
  \footnotesize\ttfamily\raggedright\setlength{\emergencystretch}{4em}
  #1\par\endgroup}
% Editorial note: what it means, or what to watch out for when writing it up.
\newcommand{\nt}[1]{\par\vspace{2pt}{\small\textit{#1}}}

\newcommand{\secnote}[1]{\vspace{-4pt}{\small #1}\par\vspace{4pt}}

\title{\vspace{-1.5cm}Findings ledger: rigid-body symmetry in flow-matching planners}
\author{PointMass3D --- SJTU internship}
\date{Compiled \today}

\begin{document}
\maketitle

\noindent
A complete inventory of what this project measured, so a manuscript can be
written from source rather than from memory. Each entry states the claim, the
number, and how it was obtained --- including the results that came out null,
the ones that stayed inconclusive, and the nine claims that turned out to be
wrong and were retracted.

\medskip\noindent
The weight tags are a suggestion, not a verdict. \HEAD{} means it can carry a
section; \SUPP{} means it earns a paragraph; \DET{} means a footnote or a
methods line; \RETR{} means it was believed and is false; \OPEN{} means it is
not finished. Disagreeing with these is the part worth practising.

\tableofcontents

% ===================================================================
\section{Benchmark and protocol}
\secnote{Not findings. This is the material a methods section is made of, and
every number below is meaningless without it.}

\finding{A1}{\DET} The task: a 3D point mass crossing a cluttered box.
\ev{workspace [-1,1]\^{}3; 20 spheres + 20 oriented boxes per environment;
robot radius 0.03; trajectories are 64 waypoints x 3 coordinates}

\finding{A2}{\DET} Dataset and split.
\ev{300 environments x 36,000 expert trajectories (600 pairs x 30, time-reversed
pairs included) from RRT-Connect + CHOMP; training draws from envs [0,N);
evaluation always on envs 250--299, held out entirely}

\finding{A3}{\DET} The model and the generative objective.
\ev{conditional flow matching, rectified/OT path, Gaussian prior; dilated
temporal ConvNet, FiLM conditioning, PointNet obstacle encoder; 2,161,283
parameters (reduced arm); boxes as centre + three half-edge vectors (12-d)}

\finding{A4}{\DET} The evaluation protocol, fixed across every arm.
\ev{50 held-out envs x 10 distinct pairs x 20 samples = 500 problems; 8 Euler
steps; K\_FA = 1 unless stated; headline metric is per-sample collision-free \%}

\finding{A5}{\SUPP} Per-sample and best-of-20 differ by about 30 points and must
never be conflated.
\ev{world frame at 20 envs: per-sample 13.6\% vs best-of-20 44.8\%}
\nt{A real trap, not bookkeeping --- the world-frame model learned sample
\emph{diversity} without learning sample \emph{quality}. Any sentence quoting
one of these numbers has to say which.}

% ===================================================================
\section{The core result}
\secnote{The reduction is a change of coordinates: translate and rotate so start
and goal land on $(\mp d/2, 0, 0)$. No parameters, no loss term, no augmentation.}

\finding{B1}{\HEAD} At convergence, 60 environments, the reduction is worth
$+36.5$ points.

\begin{center}\small
\begin{tabular}{lrr}
\toprule
Arm & Collision-free & Seeds \\
\midrule
World frame      & $14.60 \pm 0.20$ & 3 \\
$(s,g)$ reduction & $\mathbf{51.10 \pm 0.25}$ & 3 \\
\bottomrule
\end{tabular}
\end{center}
\ev{300 epochs; errors are ACROSS SEEDS, the right floor for a claim about a
method. The spread over held-out problems answers a different question.}

\finding{B2}{\HEAD} At 250 environments the gap is larger still: $+41.4$.
\ev{converged: world frame 14.4 -> reduction 55.8}

\finding{B3}{\SUPP} The original $2\times4$ grid at a fixed 20-epoch budget,
three seeds per cell.

\begin{center}\small
\begin{tabular}{lrrrr}
\toprule
Environments & 20 & 60 & 150 & 250 \\
\midrule
World frame       & 12.86 & 13.57 & 14.75 & 15.37 \\
$(s,g)$ reduction & 19.23 & 34.43 & 40.60 & 45.45 \\
\bottomrule
\end{tabular}
\end{center}
\nt{Useful as the scaling curve; superseded by \texttt{B1}/\texttt{B2} as a
statement of what the method achieves. \textbf{Three headline levels exist and
they are not interchangeable.} $51.10\pm0.25$ against $14.60\pm0.20$ is the
headline: 60 environments, converged, three seeds, and the setting every
mechanism comparison uses. $45.5$ (this grid, at 250 environments) runs at a
fixed 20-epoch budget and answers how the gap moves with data, not how well the
method does. $55.8$ is the largest number measured, 250 environments converged,
and rests on one seed, which is why it is not the headline despite being the
largest. Label every one of them wherever it appears --- four review rounds were
spent on exactly this.}

\finding{B4}{\SUPP} Seed-matched gaps, which is the correct paired test.
\ev{20 envs +6.37; 60 envs +20.9 +- 1.6; 150 envs +25.85; 250 envs +30.08 +- 0.38}

% ===================================================================
\section{How it scales}

\finding{C1}{\HEAD} The world-frame arm is flat in data: $+2.6$ points across a
$12.5\times$ increase in layouts.
\ev{20 -> 250 envs: 12.86 -> 15.37}

\finding{C2}{\HEAD\ \OPEN} The reduction gains $+26.6$ over the same range and
has \emph{not} plateaued at 250 environments --- which is the dataset cap.
\ev{19.23 -> 45.45 over 20 -> 250 envs, with 50 envs held out}

\finding{C3}{\HEAD} The gap widens with data --- the opposite of the
``small-data crutch'' hypothesis.
\nt{Worth stating explicitly as a refutation. The natural objection to a
representation trick is that it only helps when data is scarce; here it helps
\emph{more} as data grows.}

\finding{C4}{\SUPP} Minimum clearance agrees independently of the collision-free
rate.
\ev{world frame -0.078 -> -0.070 (static); reduction -0.058 -> -0.0095
(approaching feasibility)}
\nt{Two metrics moving together is a much stronger argument than either alone.}

% ===================================================================
\section{Controls: does the result survive scrutiny}
\secnote{Each of these exists because someone objected, or would have.}

\finding{D1}{\HEAD} \textbf{Not underfitting.} The control gets nothing from a
$30\times$ budget increase.
\ev{300 epochs at 250 envs: world frame 14.6 -> 14.4; reduction 38.8 -> 55.8}
\nt{This answers the review objection directly. The control has \emph{converged
onto a worse solution}, which is a different and stronger claim than ``it needs
more training''.}

\finding{D2}{\HEAD} \textbf{Not an artefact of flow matching.} A DDPM arm
reproduces the gap.
\ev{eps-prediction, cosine schedule; at 20 epochs +16.71; at convergence +37.0
against flow matching's +36.5}
\nt{The two objectives disagree early and agree at convergence --- itself a
small finding about how much early-training comparisons are worth.}

\finding{D3}{\HEAD} At convergence the world-frame arm does not \emph{beat} the
straight-line baseline.
\ev{world frame 14.58 +- 0.22 (n=3) vs straight line 15.6, per sample. PAIRED
bootstrap over environments, differencing within each replicate: -1.02 points,
95\% CI [-3.06, +0.90], SE 1.01}
\nt{The right test is paired: both arms run the identical 500 problems and both
track per-environment clutter, so their outcomes are correlated and the
difference is far better determined than either rate. Pairing halves the standard
error, 1.01 against 2.2 for either rate alone --- and the interval still spans
zero, because the effect is only 1.02. So the defensible claim is ``does not
beat'', not ``is below''. \emph{Do not quote the across-seed spread here}: it
describes variation between training runs, not uncertainty against a fixed
baseline. See \texttt{L10}. Run it with \texttt{paired\_floor\_test.py}.}

\finding{D4}{\SUPP} Obstacle-density sweep: the gap is not a property of one
clutter level.
\ev{at 8 obstacles the control tracks the straight line (61.4); the gap peaks at
24 obstacles (+41.2)}

\finding{D5}{\SUPP} Under best-of-20 the picture changes and must be reported
separately.
\ev{best-of-20 90.4 for the reduction; the world frame LOSES 3.1 points of
best-of-20 with 3x budget (39.2)}

% ===================================================================
\section{Attribution: where the benefit comes from}

\finding{E1}{\HEAD} The canonicalisation does essentially all the work; roll
augmentation adds almost nothing.
\ev{world frame 17.6 -> reduction without roll augmentation 34.1 -> with it
34.6; seed-matched the augmentation is +0.8}

\finding{E2}{\HEAD} Linear probes locate the mechanism: the reduction moves the
query$\times$scene interaction out of the network and into the representation.

\begin{center}\small
\begin{tabular}{lrr}
\toprule
Scene code & Within-env std & $R^2$ for clearance \\
\midrule
World frame   & 0.000000 & 0.015 \\
Reduced frame & 0.182    & 0.778 \\
\bottomrule
\end{tabular}
\end{center}
\nt{The world-frame scene code is \emph{exactly constant} within an environment
--- it cannot vary with the query, so it cannot encode anything about this
particular trip. The single most explanatory measurement in the project.}

\finding{E5}{\HEAD} The benefit is \emph{not} simply letting the encoder see the
query. A world-frame arm with a query-dependent scene code gains one point of the
thirty-six.

\begin{center}\small
\begin{tabular}{lrr}
\toprule
60 envs, converged, $n=3$ & free \% & vs.\ world frame \\
\midrule
world frame                    & $14.58 \pm 0.22$ & \\
$+$ query-conditioned encoder  & $15.59 \pm 0.38$ & $+1.01$ \\
$(s,g)$ reduction              & $51.12 \pm 0.27$ & $+36.54$ \\
\bottomrule
\end{tabular}
\end{center}
\ev{the arm is the world frame throughout, no change of frame anywhere, with the
raw query concatenated to every obstacle BEFORE the pool. Scene code
within-environment std 0.022 against the baseline's 0.000000; 0.07\% more
parameters. Seed-matched difference +1.01, positive on all three seeds (+1.14,
+0.50, +1.40). Built as \texttt{--query-encoder}.}
\nt{This is the answer to the strongest objection the work faces. \texttt{E2}
shows the world-frame scene code cannot vary with the query \emph{by
construction}, which invites the deflationary reading that the reduction buys
query-dependence rather than canonicalisation. It buys about one point of the
36.5; canonicalisation supplies 35.5. \emph{Second observation, and the better
one}: at 15.59 the arm lands ON the straight line's 15.6, not above it. Query
dependence is exactly enough to reach the trivial baseline and not enough to beat
it. What takes the model past the floor is the change of frame. Related:
[[E2]], [[D3]].}

\finding{E3}{\SUPP} The five removed degrees of freedom decompose, and they
interact.
\ev{3 translations +12.3; 2 rotations +18.0; interaction +7.1; total +30.1}
\nt{Neither half explains the result alone.}

\finding{E4}{\SUPP} The reduction also makes each additional epoch more
productive.
\ev{gain from extra epochs: reduced arm +8.0; world frame +0.9}

% ===================================================================
\section{The residual symmetry, and three ways to impose it}
\secnote{Five degrees of freedom go exactly. The sixth --- a roll about the
start--goal axis --- cannot be removed, so it must be handled some other way.}

\finding{F1}{\HEAD} The reduction collapses SE(3) to SO(2), not to nothing.
\ev{verified to 1e-15 in both domains; every reduced quantity maps through one
shared rotation about the first axis; invariant outright: d, the x-coordinate of
every point, and its radius in the y-z plane}
\nt{This is what \emph{motivates} roll augmentation and frame averaging. If the
representation were fully invariant, both mechanisms would be unmotivated --- so
the correction (\texttt{L1}) made the paper more coherent, not less.}

\finding{F2}{\SUPP} The measured non-equivariance residual is small and
\emph{flat} across data scale.

\begin{center}\small
\begin{tabular}{lrrrr}
\toprule
Environments & 20 & 60 & 150 & 250 \\
\midrule
$r$ (RMS, $K=3$) & 0.0183 & 0.0170 & 0.0177 & 0.0168 \\
\bottomrule
\end{tabular}
\end{center}
\ev{no-roll checkpoint at 60 envs: 0.0279}

\finding{F3}{\SUPP} Frame averaging is small but not inert, and its value
depends on the training budget.
\ev{at 20 epochs the K sweep is flat: 34.6 / 35.1 / 35.4 / 35.2 / 35.3 for
K = 1/3/5/7/9; at convergence, 60 envs, 3 seeds: +0.68 +- 0.13; at 250 envs
converged: +1.57; cost at K=9 is 1.05x wall time}

\finding{F4}{\DET} Frame averaging improves endpoint error by about 35\%, and
that does not reach the collision-free rate.
\ev{endpoint error 0.0020 -> 0.0013, consistent on every checkpoint}
\nt{A clean example of a real effect on the wrong quantity --- endpoint
precision was never the failure mode.}

\finding{F5}{\SUPP} Frame averaging is the orthogonal projection onto the
equivariant subspace, so $r$ bounds what it can possibly remove.
\ev{symmetrisation budget at 60 envs: under 3e-4}
\nt{A theorem, and the cleanest piece of theory in the paper. Note carefully
what it does \emph{not} say --- see \texttt{F6}.}

\finding{F6}{\HEAD} The residual does not predict what a symmetry mechanism is
worth.

\begin{center}\small
\begin{tabular}{lrr}
\toprule
Arm & $r$ under SE(3) & Collision-free \\
\midrule
World frame                      & 0.0881 & 15.4 \\
World frame + SE(3) augmentation & 0.0181 & 17.5 \\
$(s,g)$ reduction                & 0.0174 & 45.6 \\
\bottomrule
\end{tabular}
\end{center}
\nt{Augmentation and canonicalisation leave the field \emph{equally
equivariant} and differ by 28 points. The tempting rule ``build the constrained
architecture only if $r$ is large'' has no support. This turned a caveat a
reviewer asked for into a measured finding --- a good story to tell.}

\finding{F7}{\SUPP} Proposition 1 verified on a \emph{trained} model, not just
in theory.
\ev{SE(3) residual 0.0186 (K=9) / 0.0194 (K=32) vs roll residual 0.0191 / 0.0192
at matched K --- agreement to 0.0005}

% ===================================================================
\section{The equivariant architecture: the third mechanism}
\secnote{The published version said this was untested, and that the projection
bound could not stand in for it. This is the work that tested it.}

\finding{G1}{\HEAD} An exactly SO(2)-equivariant backbone, verified before
training.
\ev{features split by irrep -- along-axis component invariant (m=0),
perpendicular pair rotating (m=1); complex-linear channel mixing;
invariant-weighted attention pooling over obstacles; normalisation by an
invariant; equivariance error 1.2e-06 on an UNTRAINED network}

\finding{G2}{\SUPP} The equivariance test has a tightness control, so it cannot
pass vacuously.
\ev{off-axis rotation gives error 5.03; equivariance survives multiplying every
weight by 3 (relative error 2.1e-06)}
\nt{A network that ignored its vector inputs entirely would pass the
equivariance test and fail this one. Worth a sentence --- most equivariance
claims are not checked this way.}

\finding{G3}{\HEAD} The constrained model is 19.5\% \emph{larger} than the
baseline, and still no better asymptotically.
\ev{equivariant 2,582,233 vs unconstrained 2,161,283 parameters}
\nt{Stronger than capacity-matching would have been: it rules out ``the
constrained model was simply smaller'' without any argument. See \texttt{L9}.}

\finding{G4}{\HEAD} Equivariance buys training efficiency, not a higher ceiling.

\begin{center}\small
\begin{tabular}{lrrr}
\toprule
Epoch & Equivariant & Unconstrained & $\Delta$ \\
\midrule
10  & 35.4 & 22.1 & $+13.3$ \\
40  & 44.6 & 37.9 & $+6.7$ \\
100 & 48.7 & 45.8 & $+2.9$ \\
300 & 51.27 & 51.12 & $+0.15$ \\
\bottomrule
\end{tabular}
\end{center}
\ev{Seed-matched over THREE seeds. The optimisation saving, computed by
interpolating when the unconstrained arm reaches each constrained level: 3.2x at
epoch 10, 2.4x at 20, 2.2x at 40 and 60.}
\nt{Quote \emph{+0.15} at the ceiling and \emph{two to three times sooner} for
the saving. Two earlier numbers are wrong and both came from seed 0 alone: +0.05
for the ceiling, and ``fourfold'' for the saving, which is 2.2--3.2x on
multi-seed means. See \texttt{L11}.}
\nt{All three mechanisms now agree the residual roll is nearly exhausted. That
\emph{supports} the central claim rather than leaving a hole in it.}

\finding{G5}{\SUPP} The first implementation scored 22\% --- and the failure was
informative.
\ev{v1 plateaued at 22.2 and overfitted from epoch 40 (val loss 0.0475, best
epoch ~40); v2 reaches val 0.0380 with best epoch 280}
\nt{Two defects, both invisible to the equivariance tests because the
mathematics was correct throughout: mean-pooling the obstacle direction features
retained only 14\% of a typical obstacle's magnitude (the centroid of a roughly
symmetric cloud is near zero), and the rotating stream had no normalisation
while the scalar stream had two per block. \emph{Correct and useless are
different failure modes, and only one of them has a test.}}

\finding{G6}{\SUPP} The efficiency advantage persists with more data.
\ev{250 envs: equivariant 54.8 at epoch 80 vs unconstrained 51.6 -- it reaches
by epoch 80 what the unconstrained arm needs about epoch 200 for}

% ===================================================================
\section{Local geometry: the unplanned result}
\secnote{Built as a diagnostic to measure how much headroom a better obstacle
encoder could reach. It turned out to be the largest single effect in the
project.}

\finding{H1}{\HEAD} What it is --- and a distinction the paper must get right.
\ev{at each waypoint, the signed distance to the scene and its gradient; split
by irrep for the equivariant arm (d and g\_x invariant, g\_yz rotating)}
\nt{This is \emph{not privileged information}: it is a deterministic function of
the obstacle primitives the model already receives, so it is a featurisation,
not an oracle. The real scope limit is narrower --- it assumes obstacles are
given as explicit primitives at inference, true here and false if you plan from
images. State that limit yourself.}

\finding{H2}{\HEAD} The $2\times2$.

\begin{center}\small
\begin{tabular}{lrrrr}
\toprule
60 envs, 300 epochs & No local geom. & $+$ local geom. & Geometry adds & $n$ \\
\midrule
World frame               & $14.60 \pm 0.22$ & $54.54 \pm 0.01$ & $+40.0$ & 3 / 2 \\
$(s,g)$ reduction         & $51.10 \pm 0.27$ & $79.37 \pm 0.10$ & $+28.3$ & 3 / 3 \\
$+$ equivariant backbone  & $51.27 \pm 0.22$ & $82.13 \pm 0.02$ & $+30.9$ & 3 / 2 \\
\midrule
\emph{reduction adds}     & $+36.5$ & $+24.8$ & & \\
\bottomrule
\end{tabular}
\end{center}
\ev{every cell now n>=2; the local-geometry arms are the least seed-sensitive
measurements in the project, SE 0.01 and 0.10 against 0.22 and 0.27 without}

\finding{H3}{\HEAD} The reduction survives a good encoder: still worth $+24.7$
with local geometry present.
\ev{reduction alone +36.2; reduction with local geometry +24.7}
\nt{This defuses the obvious objection --- that the reduction was only
compensating for a weak obstacle encoder. It was run specifically to find out
whether it sank the earlier result. Say that.}

\finding{H4}{\HEAD} Local geometry alone ($+39.9$) is a bigger lever than the
reduction alone ($+36.2$).
\nt{The mechanism is no longer the largest effect on its own benchmark. It is
still large, still free, and still composes --- but ``symmetry is the main
lever'' is \emph{no longer supportable}. A paper that reports this itself reads
as careful; a reviewer who finds it reads as catching you.}

\finding{H7}{\HEAD} Local geometry scales better in data than the query frame
does, and the gap between them widens.

\begin{center}\small
\begin{tabular}{lrrr}
\toprule
Budget-matched & world frame & $+$ reduction & $+$ local geom. \\
\midrule
60 envs, epoch 300  & 14.60 & $51.10$ ($+36.5$) & $54.54$ ($+40.0$) \\
250 envs, epoch 80  & 14.5  & $51.6$ ($+37.1$)  & $63.4$ ($+48.9$) \\
\bottomrule
\end{tabular}
\end{center}
\ev{The decay, seed-matched and budget-matched at 80 epochs, two seeds at each
scale: the reduction alongside local geometry is worth +22.37 +- 0.46 at 60
environments and +17.24 +- 0.17 at 250. A 5.1-point decay against SEs of 0.46 and
0.17, so it is measured, not suggested.}
\nt{The reduction is flat across a 4x increase in environments at matched budget;
local geometry gains 11 points. Sharpest form: the world-frame arm with local
geometry reaches \emph{57.3 at epoch 20}, beating the canonicalised arm without
it at epoch 300 (55.8). The mechanism is consistent with [[E2]]: the global
encoder pools 40 obstacles into one vector applied identically to every waypoint,
so more environments means more scenes through the same bottleneck, while local
geometry bypasses it. \emph{The first version of this comparison was confounded}
--- it read 60 envs at epoch 300 against 250 envs at epoch 80, and about a third
of the apparent decay was budget rather than data.}

\finding{H5}{\SUPP} The two effects are sub-additive.
\ev{additive prediction 14.6 + 36.2 + 39.9 = 90.7; measured 79.2 -- short by 11.5}
\nt{Expected and honest: both attack the same underlying difficulty (relating
query geometry to scene geometry), and neither subsumes the other.}

\finding{H6}{\SUPP} Equivariance and local geometry interact: $+0.1$ alone,
$+3.0$ together.
\ev{without geometry 50.8 -> 50.9; with geometry 79.2 -> 82.2}
\nt{Mechanistically what you would predict: the rotating stream only pays off
once per-waypoint directional information exists for it to act on.}

% ===================================================================
\section{Baselines and calibration}
\secnote{All on the same 500 held-out problems, same collision checker, same
pair selection. This is what makes the learned numbers interpretable.}

\finding{I1}{\HEAD} Classical reference points.

\begin{center}\small
\begin{tabular}{lrrr}
\toprule
Planner & Free \% & Clearance & s/query \\
\midrule
Straight line        & 15.6  & 0.055 & $<$0.001 \\
TrajOpt              & 74.0  & 0.026 & 0.352 \\
CHOMP                & 88.8  & 0.019 & 0.394 \\
RRT-Connect          & 99.6  & 0.017 & 0.284 \\
Expert (RRT+CHOMP)   & 100.0 & 0.037 & 0.342 \\
Learned planner      & ---   & ---   & 0.067 \\
\bottomrule
\end{tabular}
\end{center}
\nt{Single-core CPU for the classical planners, one GPU for the learned one, so
the speed comparison needs a caveat. The honest framing is \emph{4--5x faster,
far less accurate}.}

\finding{I2}{\SUPP} The untrained-prior floor, swept over width, never exceeds
15.8 best-of-20.

\begin{center}\small
\begin{tabular}{lrrrr}
\toprule
$\sigma$ & 0 & 0.05 & 0.10 & 0.30 \\
\midrule
Per sample  & 15.6 & 11.5 & 7.4  & 0.4 \\
Best-of-20  & 15.6 & 15.8 & 13.4 & 2.2 \\
\bottomrule
\end{tabular}
\end{center}
\nt{The correct control for a best-of-20 number. Widening the prior buys nothing
at any width, so the world-frame arm's 44.8\% best-of-20 is real diversity.}

\finding{I3}{\DET} Clustered bootstrap over environments confirms the assumed
standard errors.
\ev{10,000 resamples; SE 2.2pp at the 15.6\% rate, 2.9pp at 74.0\%}

% ===================================================================
\section{Structural and theoretical results}

\finding{J1}{\HEAD} Canonicalising the sixth degree of freedom is impossible, by
the hairy ball theorem.
\ev{a gauge y(x) perpendicular to x is exactly a unit tangent field on the
2-sphere; x covers the whole sphere here, so there is no contractible-domain
escape; every deterministic gauge jumps somewhere -- for the smallest-component
cross-product construction the jump is 60--120 degrees}

\finding{J2}{\HEAD} A scene-derived gauge fails too, by central inversion.
\ev{obstacle centres uniform on a cube are symmetric under p -> -p, which
preserves every radius and sends phi -> phi + pi, so the expected m-th angular
moment vanishes for every ODD m and every axis direction}
\nt{The gauge anyone would actually build --- ``point $\hat{y}$ at the obstacle
mass'' --- \emph{is} the $m=1$ moment, so this kills it. Even moments survive
this but not \texttt{J1}: an $m=2$ moment defines a director, and a continuous
director field on $S^2$ is obstructed for the same Euler-characteristic reason.
See \texttt{L2} for the argument this replaced.}

\finding{J4}{\HEAD} A precise scope condition: the group must act on the STATE
space, not merely on the workspace.
\ev{holds for a point mass (state is a position) and for the SE(3) rigid body
(state is a pose); fails for a manipulator planning in joint space}
\nt{The query still defines a frame in an arm's workspace, but a rigid motion of
the world induces no clean action on joint coordinates, so there is nothing for
the reduction to apply to. Mobile bases, free-flying bodies and
end-effector-space planning pass the test; joint-space planning for arms does
not. \emph{This is why a manipulator benchmark would likely demonstrate that the
method does not apply rather than that it does.} Give a reader the test rather
than a third domain on which it happens to work.}

\finding{J3}{\SUPP} A second domain exists: SE(3) rigid bodies, with the same
qualitative result.
\ev{L-shaped union of 5 spheres (trivial symmetry group, deliberately -- a rod
would make the domain secretly SE(3)/SO(2)); 9-d state, position + 6-d rotation;
conditioning 18 -> 13 under the reduction, not 6 -> 1; world frame 3.0 ->
reduction 8.6, trivial floor 4.8, ratio 2.9x}
\nt{This answers ``a point mass in $\mathbb{R}^3$ is the easy case'' --- the
strongest single objection to the benchmark.}

% ===================================================================
\section{Null and inconclusive results}
\secnote{Worth writing up. A paper that reports only what worked is less
believable than one that reports what did not.}

\finding{K1}{\DET} The cone experiment failed to decouple residual from quality.
\ev{r rises 31\% as the cone narrows (0.0238 at 15 degrees) but quality falls
with it; cone-matched re-evaluation lifted all arms about 15 points without
equalising quality (12-point spread); among the three arms whose quality does
match, r barely varies}
\nt{Report as \emph{an attempted decoupling that failed}, not as evidence either
way.}

\finding{K2}{\SUPP} Roll augmentation is not what makes the model
roll-equivariant.
\ev{no-roll residual 0.0279 vs 0.0170 with augmentation -- higher, both tiny}
\nt{The frame construction is a deterministic function of the axis direction,
and 60 envs x 600 pairs supply that many distinct directions, so the
reduced-frame data already spans a wide spread of rolls. \emph{Explicit
augmentation solves an already-solved problem.} A prediction that turned out
wrong.}

\finding{K3}{\DET} Frame averaging on the no-roll checkpoint drifts slightly
negative.
\ev{34.1 / 33.6 / 33.2 / 33.2 for K = 1/3/5/9; within noise (~4pp SE), but
clearance worsens while length shortens -- the signature of straighter paths
cutting nearer obstacles}
\nt{Flag, do not conclude.}

\finding{K4}{\DET} The $K$ sweep for frame averaging is flat from $K=3$ upward,
and cheap at any $K$.
\ev{no gain past K=3 at any budget; K=9 costs 1.05x wall time}

% ===================================================================
\section{Retracted claims}
\secnote{Every one of these appeared in a draft and was wrong. They are listed
because the corrected version is usually the more interesting sentence, and
because writing one of them again by accident is the easiest mistake available.}

\finding{L1}{\RETR} ``The reduced representation is bit-identical under any
rigid transformation.''
\nt{False. It is invariant only up to a shared roll --- see \texttt{F1}. The
frame's tie-break axis is a fixed world axis, so rotating the problem generally
selects a different one. The correction is what motivates the whole second half
of the paper.}

\finding{L2}{\RETR} ``Cube symmetry is $C_4$ in every coordinate plane, so all
moments below $m=4$ vanish.''
\nt{Holds only when the axis is a four-fold axis. Measured $|E[c_2]| = 0.070$ on
a face diagonal against a 0.002 noise floor. Replaced by central inversion
(\texttt{J2}), which is simpler and stronger.}

\finding{L3}{\RETR} ``The Brownian-bridge prior is not isotropic in the plane
normal to the axis.''
\nt{False --- an i.i.d.\ bridge has covariance $\sigma^2 g(1-g)I$ and is
roll-invariant. The real costs are query-dependence and the singular endpoint
covariance.}

\finding{L4}{\RETR} ``A small residual tells you in advance not to build a
constrained architecture.''
\nt{Refuted twice: by \texttt{F6} (equal residuals, 28-point gap) and then by
\texttt{G4} (the constrained architecture \emph{is} worth building, for
efficiency). The published draft used this to justify not building it. That
reasoning was wrong.}

\finding{L5}{\RETR} ``The residual falls with training data.''
\nt{Came from two points under an older mean-of-norms definition. Under the RMS
definition across four scales it is \emph{flat}: 0.0183 / 0.0170 / 0.0177 /
0.0168.}

\finding{L6}{\RETR} ``Frame averaging and roll augmentation are inert.''
\nt{True at 20 epochs, false at convergence: $+0.68 \pm 0.13$ at 60 envs,
$+1.57$ at 250. Only the $K$ sweep is genuinely flat.}

\finding{L7}{\RETR} ``The world-frame arm is indistinguishable from joining the
endpoints.''
\nt{At convergence it is \emph{below} the straight line --- a sharper claim.}

\finding{L8}{\RETR} ``Roll augmentation is worth $+4.4$.''
\nt{An unpaired artefact. Seed-matched it is $+0.8$.}

\finding{L10}{\RETR} ``The world-frame arm is roughly five across-seed standard
errors below the straight line.''
\nt{Wrong denominator. The across-seed spread describes variation between
training runs, not the uncertainty of a comparison against a fixed deterministic
baseline measured on the same problems. The right test is paired
(\texttt{D3}): SE 1.01, interval [-3.06, +0.90], which includes zero. Note the
correction cuts both ways --- pairing HALVED the standard error, so the original
error bars were too wide as well as attached to the wrong quantity, and the
effect is simply too small either way.}

\finding{L11}{\RETR} ``The equivariant backbone gives a fourfold saving in
optimisation, and $+0.05$ at the ceiling.''
\nt{Both came from seed 0 alone. On multi-seed means the saving is 2.2--3.2x
depending where it is measured, and the ceiling difference is +0.15. Third
instance of the same failure: a seed-0 figure that became an overstatement once
seeds landed. The first was roll augmentation at +4.4, which is +0.8 paired
(\texttt{L8}); the second was the query-encoder control at 16.15, which is 15.59
over three seeds. \emph{Rule, now applied: never quote a single-seed difference
against a multi-seed mean.}}

\finding{L9}{\RETR} ``The equivariant model's capacity is matched to the
baseline within 0.8\%.''
\nt{It is 19.5\% \emph{larger} --- 2,582,233 against 2,161,283. The module
defaults were matched; the training script overrode them from shared CLI flags
and nothing checked. See \texttt{G3}: the true statement is stronger than the
false one.}

% ===================================================================
\section{Method notes worth a paragraph each}
\secnote{Not results, but the kind of detail that makes a methods section
trustworthy, and that a reader can act on.}

\finding{M1}{\SUPP} 61 automated tests across six files, all property-based.
\ev{geometry 10; frame averaging 9; diffusion schedule 7; SE(3) 12; SDF 13;
equivariant backbone 10; every equivariance test runs on an UNTRAINED network}

\finding{M2}{\SUPP} Points versus free vectors is a silent-failure class.
\nt{Waypoints, start, goal and obstacle centres take the full affine map; the
three box half-edge vectors take the \emph{rotation only}. Getting it wrong
destroys box extent while preserving box position --- the model goes
obstacle-size-blind with a perfectly clean loss curve. The same class bit the
SE(3) domain's rotation columns.}

\finding{M3}{\SUPP} Two estimator traps in the residual probe.
\nt{The probe must draw where the model's own sampler draws: drawing in world
coordinates for the reduced arm inflated $r$ by $3.4\times$ (0.061 vs 0.017).
And an identity-decode test cannot catch it --- identity is equivariant under
all rotations, so both branches return zero. The test that catches it is a field
equivariant to \emph{rolls only}.}

\finding{M4}{\DET} A differentiable batched SDF, cross-validated against the
simulator.
\ev{agrees with the collision checker to 8e-8; rigid-motion equivariant to 2e-7}

\finding{M7}{\SUPP} Three tools now guard the numbers, and each exists because
something got past the previous one.
\nt{\texttt{check\_paper\_numbers.py} asks ``does this value appear, and has a
retired one come back''. It cannot catch a \emph{claim about} a number: a caption
saying $n{=}2$ while a paragraph six pages away says single-seed contains no
value it knows. \texttt{seed\_inventory.py} generates the seed count per arm from
\texttt{results/conv/}, takes the last epoch PER SEED, and refuses to average
across mixed budgets --- a killed run once contributed an epoch-20 score as if it
were a seed. \texttt{audit\_numbers.py} prints every occurrence of every headline
value with context, because the recurring failure was never a wrong number but a
right number changed in one place and not the three others it lives in. Run the
third after every change; it is the loop that was being skipped.}

\finding{M8}{\SUPP} Two ways a results table looked plausible and was wrong.
\nt{First, \texttt{*.ep0*.json} also matches \texttt{ep0300-kfa3.json}, so a
summary taking the last match per seed silently reported frame-averaged scores
for whichever arm had a K sweep --- the treatment read 51.64 instead of 51.12,
inflated by exactly the frame-averaging gain. Second, a filter written
\texttt{if (envs=='250' and ep!=80) and ep!=300} reduces to \texttt{False} for
every 60-environment path, so no filtering happened at all and each cell showed
whatever epoch came last in glob order. \emph{Always print the epoch a row came
from.}}

\finding{M5}{\DET} The manuscripts' numbers are checked mechanically, not by
reading.
\nt{\texttt{check\_paper\_numbers.py} holds about 50 values that must appear,
about 25 that must \emph{not} (superseded ones), and a set that must match the
figure-generation constants --- which drift independently of the prose. That
last category caught a scaling figure plotting single-seed values beside a
three-seed table.}

\finding{M6b}{\DET} Two globbing traps that silently corrupted a results table.
\nt{First: \texttt{*.ep0*.json} matches \texttt{ep0300-kfa3.json} as well as
\texttt{ep0300.json}, so a summary that takes the last match per seed silently
reports frame-averaged scores for whichever arm had a K sweep run on it --- the
treatment arm read 51.64 instead of 51.12, inflated by exactly the frame
averaging gain. Second: a filter written as
\texttt{if (envs=='250' and ep!=80) and ep!=300} reduces to \texttt{False} for
every 60-environment file, so no filtering happens at all. Both produce a table
that looks entirely plausible. Always print the epoch a row came from.}

\finding{M6}{\DET} Seed-matched comparisons are only valid within a device.
\ev{torch.Generator draws differ between CPU and CUDA for the same seed --- local
eval gave 13.6\% against 12.8\% reported from the GPU run}

% ===================================================================
\section{Open, and honestly unfinished}

\finding{N1}{\OPEN} Seed inventory as it stands.
\ev{n=3: ctrl, treat, equiv2, treatgeo (60 envs), ctrlqc. n=2: ctrlgeo, equivgeo
(60 envs), ctrlgeo, treatgeo (250 envs). n=1: ctrl, treat, equiv2 (250 envs).}
\nt{Every claim the paper leans on has at least two seeds. The three single-seed
cells are the plain 250-environment arms, which support the scaling narrative and
are labelled as point estimates. Regenerate with \texttt{seed\_inventory.py}
rather than quoting this from memory.}

\finding{N2}{\OPEN} One architecture, one clutter generator, one workspace
family.
\nt{The cheapest of the three to attack is capacity: doubling the trunk
(\texttt{--channels 256}) on both arms at 60 environments, three seeds, is six
runs of about four hours. If the gap survives at 4x capacity then ``the model is
small'' stops being a limitation and becomes a controlled variable. The other two
need a different benchmark, and \texttt{J4} explains why the obvious candidate
does not qualify.}

\finding{N3}{\OPEN} The reduction has not plateaued at 250 environments, which
is the dataset cap.
\nt{Cutting trajectories-per-pair from 30 to 5 would buy about $6\times$ the
environments for the same generation time, and environments are what the curve
says are binding.}

\finding{N4}{\OPEN} Whether local geometry survives outside explicit-primitive
obstacle representations is untested.
\nt{The honest scope sentence for \texttt{H1}, and a natural future-work
paragraph.}

% ===================================================================
\section*{A suggestion about shape}
\addcontentsline{toc}{section}{A suggestion about shape}

If you want one, the spine that holds all of this together is: \emph{a
query-derived frame removes five degrees of freedom exactly and for free; the
sixth is provably not removable; and all three ways of handling the sixth agree
that it barely matters.} That is a complete argument with a theorem
(\texttt{J1}, \texttt{J2}), a measurement (\texttt{B1}), and a triangulation
(\texttt{E1}, \texttt{F3}, \texttt{G4}).

\medskip\noindent
Section 8 then sits outside that spine as a second axis that composes with it
--- including \texttt{H4}, which says the mechanism is not the largest effect on
its own benchmark. Deciding where to put that sentence is the single most
consequential writing choice in this project. Putting it early reads as
confident.

\medskip\noindent
The material most papers would leave out and this one can afford to keep:
\texttt{F6} (a measurement that refutes an intuition the field holds),
\texttt{G5} (a bug invisible to a correct proof), and all of Section 12.
Reviewers trust a paper that argues against itself in public.

\end{document}
